\documentclass[11pt]{article}

\usepackage[T1]{fontenc}
\usepackage{lmodern}
\usepackage{amsmath,amssymb,amsthm,mathtools}
\usepackage[margin=1.08in]{geometry}
\usepackage{microtype}
\usepackage{enumitem}
\usepackage{xcolor}
\usepackage[colorlinks=true,linkcolor=blue!55!black,citecolor=blue!55!black,
  urlcolor=blue!55!black]{hyperref}
\emergencystretch=2em

\newtheorem{theorem}{Theorem}[section]
\newtheorem{proposition}[theorem]{Proposition}
\newtheorem{lemma}[theorem]{Lemma}
\newtheorem{corollary}[theorem]{Corollary}
\theoremstyle{definition}
\newtheorem{definition}[theorem]{Definition}
\theoremstyle{remark}
\newtheorem{remark}[theorem]{Remark}

\newcommand{\C}{\mathbb C}
\newcommand{\N}{\mathbb N}
\newcommand{\cA}{\mathcal A}
\newcommand{\cB}{\mathcal B}
\newcommand{\cD}{\mathcal D}
\newcommand{\cE}{\mathcal E}
\newcommand{\cK}{\mathcal K}
\newcommand{\cM}{\mathcal M}
\newcommand{\cP}{\mathcal P}
\newcommand{\cQ}{\mathcal Q}
\newcommand{\Sat}{\operatorname{Sat}}
\newcommand{\Sym}{\operatorname{Sym}}
\newcommand{\dtrdeg}{\operatorname{dtrdeg}}
\newcommand{\diff}{\partial}

\title{Prime-Clock Linearization of Differential-Algebraic Closure\\
and the Euler-Completion Barrier}
\author{Michael Haufschild \and GPT 5.6 Sol \and Fable 5 \and Grok 4.5}
\date{}

\begin{document}
\maketitle

\begin{abstract}
We identify the differential-algebraic closure geometry of a universal class of
prime-power Dirichlet series.  Let
\[
  F_a(s)=\sum_p\sum_{r\geq1}a_{p,r}p^{-rs}
\]
have finite abscissa, quotient its coefficient array by arrays supported on finitely many
prime rays, and let \(X\in\C[X]\) act through the diagonal clock
\((p,r)\mapsto r\log p\).  The resulting prime-clock module \(\cQ\) linearizes the entire
relative differential-algebraic geometry of these carriers.

For every clock submodule \(M\leq\cQ\), the pullback of Kolchin differential-algebraic
closure over the carrier field \(\Phi(M)\) is exactly polynomial saturation:
\[
  F_a\text{ is differential-algebraic over }\Phi(M)
  \quad\Longleftrightarrow\quad
  [a]\in\Sat_{\C[X]}(M).
\]
Simultaneous differential independence is injectivity of the corresponding free
\(\C[X]\)-module map into \(\cQ/M\), and every differential relation is generated by the
associated affine clock-linear relations and their prolongations.  Thus a nonlinear closure
problem is represented by linear module geometry.

Cofinally recurrent replication embeds a profile space \(E\) as the free saturated sector
\(\C[X]\otimes_{\C}E\).  Consequently analytic completion and differential-algebraic closure
are maximally transverse: a replicated profile is differential-algebraic over the field of a
subspace \(W\leq E\) exactly when it already belongs to \(W\).  For scalar valuation profiles,
the finite-support field is
\(\cB_{\rm loc}\langle P(ms):m\geq1\rangle_{\diff}\), whereas the Euler profile is
\((1/m)_{m\geq1}\).  Hence, on \(\Re s>1\),
\[
  \log\zeta(s)\text{ is hypertranscendental over }
  \cB_{\rm loc}\langle P(ms):m\geq1\rangle_{\diff},
\]
although \(\log\zeta(s)=\sum_{m\geq1}P(ms)/m\) there.
\end{abstract}

\noindent\textbf{Keywords.}
Hypertranscendence; Dirichlet series; differential algebra; prime zeta function;
Euler product; module saturation; analytic completion.

\medskip
\noindent\textbf{2020 Mathematics Subject Classification.}
Primary 11J91, 12H05; Secondary 11M41, 11M06.

\section{Introduction}

Differential algebra is finitary: a differential field is closed under field operations and
differentiation, but not under convergent infinite summation.  Euler's identity
\begin{equation}\label{eq:euler-intro}
  \log\zeta(s)=\sum_{m\geq1}\frac{P(ms)}{m},\qquad \Re s>1,
\end{equation}
where \(P(s)=\sum_p p^{-s}\), makes this distinction concrete.  Every summand belongs to the
differential field generated by the countable family \((P(ms))_{m\geq1}\).  Does adjoining
every layer nevertheless make the completed Euler logarithm differential-algebraic?

We answer a more general question.  A prime-power coefficient array carries a canonical
diagonal clock
\[
  (p,r)\longmapsto r\log p,
\]
because differentiation multiplies its \((p,r)\)-coordinate by \(-r\log p\).  After finite
prime-ray information is placed in the coefficient field, differentiation therefore turns the
source quotient into a module over \(\C[X]\).  Our main theorem states that this module is not
merely an aid for proving independence: it represents the differential-algebraic closure
operator itself.

The representation is exact in three mutually reinforcing senses.  First, a carrier is in a
generated field precisely when its source class belongs to the generating clock submodule.
Second, differential-algebraic closure pulls back to polynomial saturation of that submodule.
Third, the full relation algebra is the linear presentation dictated by the clock module; there
are no hidden nonlinear relations.  These are not separate phenomena.  Together they say that
the differential-algebraic geometry of prime-power carriers has been linearized.

The completion consequence follows from one further geometric observation.  If profiles are
sampled recurrently along the primes, their clock module is canonically free:
\[
  \C[X]\otimes_{\C}E\hookrightarrow\cQ.
\]
Free submodules are saturated in this sector.  Thus a completion point outside a profile
subspace is not merely absent from its carrier field; it is hypertranscendental over that field.
For scalar profiles, finite support produces the prime-zeta dilates \(P(ms)\), while the
infinite profile \((1/m)\) produces \(\log\zeta\).  Formula \eqref{eq:euler-intro} is therefore
an exact example in which analytic completion crosses a differential-algebraic boundary.

Classical work of Ostrowski, Shapiro--Sparer, Amou--Katsurada, Pong, and others supplies powerful
one-source and operator-family transcendence criteria for Dirichlet series
\cite{Ostrowski,ShapiroSparer,AmouKatsurada,Pong}.  The point here is relative and
classificatory: an arbitrary already generated carrier field is recovered from its clock
submodule, its differential-algebraic closure is computed by saturation, and simultaneous
families receive a complete relation presentation.  The Euler result uses this classification
over one countably generated field; it is not obtained by passing independence through a limit.

Section~\ref{sec:carrier} constructs the localized carrier and clock module.
Section~\ref{sec:static} proves the static tail theorem.  Section~\ref{sec:master} proves the
linearization theorem.  Section~\ref{sec:recurrent} identifies the recurrent free sector, and
Section~\ref{sec:euler} derives the Euler-completion barrier.  The final section records the
precise novelty and scope.

\section{Localized carriers and the prime clock}\label{sec:carrier}

Let \(\cP\) be the rational primes and put \(\Omega=\cP\times\N_{>0}\).  Let \(\cA\) be the
complex vector space of arrays \(a=(a_{p,r})_{(p,r)\in\Omega}\) satisfying
\begin{equation}\label{eq:abscissa}
  \sum_p\sum_{r\geq1}|a_{p,r}|p^{-r\sigma}<\infty
\end{equation}
for at least one real \(\sigma\).  The associated right-end analytic germ is
\begin{equation}\label{eq:carrier}
  F_a(s)=\sum_p\sum_{r\geq1}a_{p,r}p^{-rs}.
\end{equation}
Write \(\cA_{\rm loc}\) for the arrays supported on finitely many prime rays, and let
\(\cB_{\rm loc}\) be the differential subfield of right-end meromorphic germs generated over
\(\C(s)\) by \(F_b\) for \(b\in\cA_{\rm loc}\).  This intrinsic local field is larger than a
rational field in finitely many exponentials: it contains every finite-abscissa local carrier on
each finite set of primes.

Define
\[
  \cQ=\cA/\cA_{\rm loc}
\]
and the diagonal operator \(T:\cA\to\cA\) by
\begin{equation}\label{eq:clock}
  (Ta)_{p,r}=r\log(p)a_{p,r}.
\end{equation}
Polynomial growth in \(r\log p\) can be absorbed by shifting the abscissa in
\eqref{eq:abscissa}; hence \(T\) preserves \(\cA\) and \(\cA_{\rm loc}\).  We make \(\cQ\)
a module over \(R=\C[X]\) by setting \(X[a]=[Ta]\).  Termwise differentiation gives
\begin{equation}\label{eq:clock-derivative}
  \diff^kF_a=(-1)^kF_{T^ka}.
\end{equation}

For an \(R\)-submodule \(M\leq\cQ\), define
\begin{equation}\label{eq:Phi}
  \Phi(M)=\cB_{\rm loc}\langle F_a:[a]\in M\rangle_{\diff}.
\end{equation}
Changing a representative changes its carrier by an element of \(\cB_{\rm loc}\), so this is
well defined.  Put
\begin{equation}\label{eq:saturation}
  \Sat_R(M)=\{q\in\cQ:\text{some }0\neq A\in R\text{ satisfies }Aq\in M\}.
\end{equation}

\begin{proposition}[Faithful coordinate realization]\label{prop:faithful}
The fraction field of finite-abscissa Dirichlet series embeds into the field of right-end
meromorphic germs.  The coordinate germ \(s\) is transcendental over its image.  Consequently,
algebraic relations over the actual coordinate field \(\C(s)\) lift faithfully to formal
finite-abscissa Dirichlet-series relations.
\end{proposition}

\begin{proof}
Uniqueness of absolutely convergent Dirichlet series gives injectivity on the convolution algebra
and hence on its fraction field.  Suppose, after clearing denominators, that
\[
  \sum_{j=0}^{d}s^j\sum_{n\geq1}c_{j,n}n^{-s}=0
\]
on a right half-plane.  Choose the least \(n_0\) for which
\(Q_{n_0}(s)=\sum_jc_{j,n_0}s^j\) is nonzero.  On the positive real axis, multiply by
\(n_0^s\).  Absolute convergence on a fixed line bounds the remaining tail by a polynomial in
\(s\) times \((n_0/(n_0+1))^s\).  Thus the nonzero polynomial \(Q_{n_0}(s)\) is exponentially
small as \(s\to+\infty\), a contradiction.  A second clearing of denominators proves the final
assertion.
\end{proof}

\begin{lemma}[Clock torsion-freeness]\label{lem:torsion-free}
The \(R\)-module \(\cQ\) is torsion-free.
\end{lemma}

\begin{proof}
Let \(0\neq A\in\C[X]\) and suppose \(A(T)a\) is local.  Outside finitely many prime rays,
\[
  A(r\log p)a_{p,r}=0.
\]
The clock spectrum is simple: \(r\log p=t\log q\) implies \(p^r=q^t\), hence \(p=q\) and
\(r=t\).  Since \(A\) has only finitely many roots, only finitely many coordinates can survive.
Thus \(a\) is local.
\end{proof}

\section{The static tail theorem}\label{sec:static}

The algebraic input is the exact linear matroid of the localized source space.

\begin{theorem}[Static tail matroid]\label{thm:static}
For an arbitrary family \((a_i)_{i\in I}\subset\cA\), the carriers \((F_{a_i})_{i\in I}\)
are algebraically independent over \(\cB_{\rm loc}\) if and only if the classes
\(([a_i])_{i\in I}\) are \(\C\)-linearly independent in \(\cQ\).
\end{theorem}

\begin{proof}
Both conditions have finite character, so take \(I\) finite.  A nonzero linear combination in
\(\cA_{\rm loc}\) gives a linear relation over \(\cB_{\rm loc}\).

Conversely, let the source classes be independent.  The row vectors
\[
  v(p,r)=(a_i(p,r))_{i\in I}
\]
span \(\C^I\) after deletion of any finite set of prime rays; otherwise a nonzero functional
would give a local source combination.  Successively deleting an entire selected ray yields
distinct primes \(p_i\) and exponents \(r_i\) for which \((a_j(p_i,r_i))_{i,j}\) is invertible.

Any proposed algebraic relation has finitely many coefficients.  After faithful lifting by
Proposition~\ref{prop:faithful}, those coefficients involve only finitely many local prime rays;
choose the packet outside them.  For the prime derivation
\((\partial_p f)(n)=v_p(np)f(np)\), the coefficient of
\(\prod_i p_i^{r_i-1}\) in the Jacobian determinant is
\[
  \left(\prod_i r_i\right)\det(a_j(p_i,r_i))_{i,j}\neq0.
\]
Distinctness of the prime rays prevents cancellation.  The Shapiro--Sparer Jacobian criterion
\cite[Theorem~3.1]{ShapiroSparer} gives algebraic independence over the common kernel, which
contains the selected coefficients.  Faithful realization returns the conclusion over
\(\cB_{\rm loc}\).
\end{proof}

\section{Prime-clock linearization}\label{sec:master}

For \(q=(q_i)_{i\in I}\subset\cQ\) and \(M\leq_R\cQ\), write
\begin{equation}\label{eq:lambda}
  \Lambda_{q,M}:\bigoplus_{i\in I}R\longrightarrow\cQ/M,
  \qquad (A_i)_i\longmapsto\sum_iA_iq_i+M.
\end{equation}
Let \(V_I=\C^{(I\times\N)}\) be the vector space of finite jet combinations.  Sending the
formal \((i,k)\)-jet to \(X^kq_i+M\) defines a linear map
\(\theta_{q,M}:V_I\to\cQ/M\).

\begin{theorem}[Prime-clock linearization theorem]\label{thm:master}
Let \(M\leq_R\cQ\), let \(a_i\in\cA\), and put \(q_i=[a_i]\).  Then the carrier map realizes
the following exact representation of differential-algebraic geometry.
\begin{enumerate}[label=\textup{(\roman*)},leftmargin=2.5em]
  \item The family \((F_{a_i})_{i\in I}\) is differentially algebraically independent over
  \(\Phi(M)\) if and only if \(\Lambda_{q,M}\) is injective.

  \item For every \(a\in\cA\),
  \begin{align}
    F_a\in\Phi(M)&\quad\Longleftrightarrow\quad [a]\in M,
      \label{eq:field-membership}\\
    F_a\text{ is differential-algebraic over }\Phi(M)
      &\quad\Longleftrightarrow\quad [a]\in\Sat_R(M).
      \label{eq:dacl-saturation}
  \end{align}
  In particular, \(F_a\) is hypertranscendental over \(\Phi(M)\) precisely when its class is
  non-torsion in \(\cQ/M\).

  \item If
  \[
    \Psi(L)=\{[a]\in\cQ:F_a\in L\},
  \]
  then \(\Psi(\Phi(M))=M\).  Thus \(M\mapsto\Phi(M)\) is an order embedding of clock
  submodules into embedded differential fields.

  \item The differential relation ideal of \((F_{a_i})_i\) over \(\Phi(M)\) is generated by
  the affine clock-linear relations whose homogeneous parts lie in
  \(\ker\theta_{q,M}\), together with all their prolongations.  Equivalently, after these
  relations are imposed, the jet algebra is the symmetric algebra on
  \(\operatorname{im}\theta_{q,M}\); no additional nonlinear relations remain.
\end{enumerate}
\end{theorem}

\begin{proof}
A differential-polynomial relation uses finitely many carriers, finitely many jet orders, and
finitely many coefficients from \(\Phi(M)\).  Each coefficient is a rational expression in
finitely many old carriers and their jets.  Let \(U\subset M\) be the finite-dimensional complex
span of all old clock classes occurring in those expressions.

Assume \(\Lambda_{q,M}\) is injective.  The finitely many new clock classes occurring in the
relation are linearly independent modulo \(M\), hence modulo \(U\).  A basis of \(U\), followed
by those new classes, is linearly independent in \(\cQ\).  Theorem~\ref{thm:static} makes the
corresponding old and new carriers algebraically independent over \(\cB_{\rm loc}\).  Therefore
the selected new jets are algebraically independent over the coefficient field of the proposed
relation, a contradiction.

If \(\Lambda_{q,M}\) is not injective, some nonzero finite polynomial packet satisfies
\[
  \sum_i A_i(T)[a_i]=m+b,
\]
where \([m]\in M\) and \(b\in\cA_{\rm loc}\).  Replacing \(T\) by \(-\diff\) through
\eqref{eq:clock-derivative} gives a nonzero affine linear differential relation over \(\Phi(M)\).
This proves (i).

If \([a]\in M\), field membership is immediate.  Conversely, a field expression for \(F_a\)
uses finitely many old clock classes.  If \([a]\notin M\), it is independent modulo their span,
and Theorem~\ref{thm:static} contradicts that expression.  This proves
\eqref{eq:field-membership}.

If \(0\neq A\in R\) satisfies \(A[a]\in M\), equation
\eqref{eq:clock-derivative} supplies a nonzero affine differential equation for \(F_a\) over
\(\Phi(M)\).  If \([a]\notin\Sat_R(M)\), the singleton map
\(R\to\cQ/M\), \(A\mapsto A[a]+M\), is injective; part (i) says that \(F_a\) is
hypertranscendental.  This proves \eqref{eq:dacl-saturation}.  Part (iii) is now exactly
\eqref{eq:field-membership}.

For (iv), start with the polynomial algebra on all formal jets.  Every vector in
\(\ker\theta_{q,M}\) evaluates to an element of \(\Phi(M)\), and hence gives the stated affine
relation.  Quotienting by the ideal they generate produces the symmetric algebra on
\(\operatorname{im}\theta_{q,M}\).  Every finite linearly independent subset of this image
lifts to clock classes independent modulo \(M\); part (i) makes the corresponding carrier jets
algebraically independent over \(\Phi(M)\).  Evaluation of the symmetric algebra is therefore
injective.  Hence the linear kernel generates the whole algebraic relation ideal.  Compatibility
with \eqref{eq:clock-derivative} shows that its differential closure is generated by the same
relations and their prolongations.
\end{proof}

The central content of Theorem~\ref{thm:master} is equation
\eqref{eq:dacl-saturation}: the nonlinear Kolchin closure operator on carriers is represented by
the elementary saturation operator on a module over the principal ideal domain \(\C[X]\).
The remaining clauses certify that this representation loses neither fields nor relations.

\section{The recurrent free sector and completion}\label{sec:recurrent}

Let \(K\) be a topological space and \(x:\cP\to K\).  Call \(x\) \emph{cofinally recurrent}
if, for every \(z\in K\), every neighbourhood \(U\ni z\), and every finite set of primes
\(S\), there is \(p\notin S\) with \(x_p\in U\).

Let \(E_x\) be a complex vector space of profiles \(u=(u_r)_{r\geq1}\), with continuous
coordinates \(u_r:K\to\C\), for which
\[
  (\operatorname{Rep}_xu)_{p,r}=u_r(x_p)
\]
belongs to \(\cA\).  Write \(\mathcal R_x(u)=F_{\operatorname{Rep}_xu}\).

\begin{theorem}[Recurrent declocking]\label{thm:declocking}
If \(x\) is cofinally recurrent, then
\begin{equation}\label{eq:free-sector}
  R\otimes_{\C}E_x\longrightarrow\cQ,
  \qquad A\otimes u\longmapsto A(T)[\operatorname{Rep}_xu]
\end{equation}
is injective.
\end{theorem}

\begin{proof}
Write a nonzero finite tensor as \(\sum_jA_j\otimes u_j\), with the \(u_j\) linearly
independent.  Collecting equal powers gives \(\sum_{k=0}^{d}X^k\otimes v_k\), where
\(v_d\neq0\).  Choose \(r_0\) and \(z_0\) with \((v_d)_{r_0}(z_0)\neq0\).  On a neighbourhood
\(U\) of \(z_0\), continuity gives \(|(v_d)_{r_0}|\geq\delta>0\) and uniform bounds for all
lower \((v_k)_{r_0}\).

If the image of the tensor were local, then outside finitely many prime rays,
\[
  \sum_{k=0}^{d}(r_0\log p)^k(v_k)_{r_0}(x_p)=0.
\]
Cofinal recurrence supplies arbitrarily large such primes with \(x_p\in U\).  Divide by
\((r_0\log p)^d\).  The lower terms tend uniformly to zero, while the leading term has modulus
at least \(\delta\), a contradiction.
\end{proof}

For \(W\leq E_x\), define
\[
  \cK_x(W)=\cB_{\rm loc}\langle\mathcal R_x(w):w\in W\rangle_{\diff}.
\]

\begin{corollary}[Completion transversality]\label{cor:completion}
Let \(x\) be cofinally recurrent and \(W\leq E_x\).  Then, for every \(u\in E_x\),
\begin{equation}\label{eq:completion-law}
  \mathcal R_x(u)\text{ is differential-algebraic over }\cK_x(W)
  \quad\Longleftrightarrow\quad u\in W.
\end{equation}
More generally, a family \((\mathcal R_x(u_i))_i\) is differentially independent over
\(\cK_x(W)\) if and only if \((u_i+W)_i\) is linearly independent in \(E_x/W\).
Consequently
\[
  \dtrdeg\bigl(\cK_x(E)/\cK_x(W)\bigr)=\dim_{\C}(E/W)
\]
for every \(W\leq E\leq E_x\).
\end{corollary}

\begin{proof}
Theorem~\ref{thm:declocking} identifies the replicated sector with
\(R\otimes_{\C}E_x\).  Modulo the submodule \(R\otimes W\), it becomes
\[
  R\otimes_{\C}(E_x/W),
\]
which is torsion-free.  Polynomial saturation in this sector is therefore exactly ordinary
membership in \(W\), and free-module independence is ordinary linear independence modulo
\(W\).  Apply Theorem~\ref{thm:master} and use finite character for the rank statement.
\end{proof}

If \(W\) is proper and dense in a natural topology on \(E_x\), the analytic closure of its
carrier family may contain all completed profiles, while its differential-algebraic closure
contains exactly the carriers of \(W\).  This is the precise sense in which the two closure
operations are maximally transverse.

\section{The Euler-completion barrier}\label{sec:euler}

Take \(K\) to be a point.  An admissible scalar profile is a sequence
\(u=(u_m)_{m\geq1}\) for which
\[
  \sum_p\sum_{m\geq1}|u_m|p^{-m\sigma}<\infty
\]
for some \(\sigma\).  Its replicated carrier is
\[
  \mathcal R(u)(s)=\sum_p\sum_{m\geq1}u_mp^{-ms}.
\]
Let \(e_m\) be the profile supported at \(m\).  Then
\begin{equation}\label{eq:prime-zeta-layer}
  \mathcal R(e_m)(s)=P(ms).
\end{equation}
Let \(E_{\rm fin}\) be the finite-support profile space and put
\begin{equation}\label{eq:Kfin}
  \cK_{\rm fin}
  =\cB_{\rm loc}\langle P(ms):m\geq1\rangle_{\diff}.
\end{equation}

\begin{theorem}[Euler-completion barrier]\label{thm:euler}
The family \((P(ms))_{m\geq1}\) is jointly differentially algebraically independent over
\(\cB_{\rm loc}\).  For every admissible scalar profile \(u\),
\begin{equation}\label{eq:scalar-completion}
  \mathcal R(u)\text{ is differential-algebraic over }\cK_{\rm fin}
  \quad\Longleftrightarrow\quad u\in E_{\rm fin}.
\end{equation}
In particular,
\begin{equation}\label{eq:logzeta-ht}
  \log\zeta(s)\text{ is hypertranscendental over }
  \cB_{\rm loc}\langle P(ms):m\geq1\rangle_{\diff}
\end{equation}
on \(\Re s>1\).
\end{theorem}

\begin{proof}
The one-point sampler is cofinally recurrent.  The coordinate profiles \((e_m)_m\) are linearly
independent, so Corollary~\ref{cor:completion}, first with \(W=0\) and then with
\(W=E_{\rm fin}\), proves joint independence and \eqref{eq:scalar-completion}.

The profile \(\lambda_m=1/m\) is admissible and has infinite support.  Absolute convergence on
\(\Re s>1\) permits rearrangement:
\[
  \mathcal R(\lambda)(s)
  =\sum_p\sum_{m\geq1}\frac{p^{-ms}}m
  =-\sum_p\log(1-p^{-s})
  =\log\zeta(s),
\]
where the branch tends to zero as \(\Re s\to+\infty\).  Equation
\eqref{eq:scalar-completion} gives \eqref{eq:logzeta-ht}.
\end{proof}

The force of Theorem~\ref{thm:euler} is its countable base.  It is stronger than exclusion from
every finite window: the coefficient field already contains every prime-zeta dilation and every
finite differential-field expression involving them.  What it lacks, and what
\(\log\zeta\) supplies, is the infinite completion operation itself.

\section{Novelty, significance, and scope}\label{sec:scope}

One-source hypertranscendence of broad classes of Dirichlet series is classical
\cite{Ostrowski,AmouKatsurada,Pong,Nagoshi}; neither that fact nor the Euler identity is claimed
as new.  Shapiro--Sparer provide the Jacobian technology behind the static direction
\cite{ShapiroSparer}, while Fr\"oberg treats the classical analytic structure of the prime zeta
function \cite{Froberg}.

The new contribution is the exact relative representation theorem: the pullback of
differential-algebraic closure is polynomial saturation in one canonical prime-clock module.
Field recovery and the complete linear relation presentation certify that this is a full
linearization, rather than a sufficient independence criterion.  Recurrent declocking then
turns the representation theorem into a general completion law, and Theorem~\ref{thm:euler}
exhibits a classical analytic completion which is hypertranscendental over the field generated
by all of its individual layers.

The theorem makes no claim about zeros of \(\zeta\) and no claim toward the Riemann hypothesis.
It does impose a positive structural requirement on Euler-side approaches: a mechanism capable
of using the completed Euler logarithm must supply an operation genuinely external to finite
differential algebra---for example analytic completion, continuation, a boundary condition, or
a coupling to the archimedean functional equation.  Adjoining every finite valuation layer does
not supply that operation.

\begin{thebibliography}{99}

\bibitem{AmouKatsurada}
M.~Amou and M.~Katsurada,
\newblock Differential transcendence of a class of generalized Dirichlet series,
\newblock \emph{Illinois J. Math.} \textbf{45} (2001), no.~3, 939--948,
\newblock \href{https://doi.org/10.1215/ijm/1258138161}{doi:10.1215/ijm/1258138161}.

\bibitem{Froberg}
C.-E.~Fr\"oberg,
\newblock On the prime zeta function,
\newblock \emph{BIT} \textbf{8} (1968), 187--202.

\bibitem{Kolchin}
E.~R. Kolchin,
\newblock \emph{Differential Algebra and Algebraic Groups},
\newblock Pure and Applied Mathematics, vol.~54, Academic Press, New York, 1973.

\bibitem{Nagoshi}
H.~Nagoshi,
\newblock Hypertranscendence of \(L\)-functions for
  \(\mathrm{GL}_m(\mathbb A_{\mathbb Q})\),
\newblock \emph{Bull. Aust. Math. Soc.} \textbf{93} (2016), 388--399,
\newblock \href{https://doi.org/10.1017/S000497271500129X}{doi:10.1017/S000497271500129X}.

\bibitem{Ostrowski}
A.~Ostrowski,
\newblock \"Uber Dirichletsche Reihen und algebraische Differentialgleichungen,
\newblock \emph{Math. Z.} \textbf{8} (1920), 241--298.

\bibitem{Pong}
W.~Y. Pong,
\newblock Applications of differential algebra to algebraic independence of arithmetic
  functions,
\newblock \emph{Acta Arith.} \textbf{172} (2016), no.~2, 149--173,
\newblock \href{https://doi.org/10.4064/aa8112-12-2015}{doi:10.4064/aa8112-12-2015}.

\bibitem{ShapiroSparer}
H.~N. Shapiro and G.~H. Sparer,
\newblock On algebraic independence of Dirichlet series,
\newblock \emph{Comm. Pure Appl. Math.} \textbf{39} (1986), no.~6, 695--745,
\newblock \href{https://doi.org/10.1002/cpa.3160390602}{doi:10.1002/cpa.3160390602}.

\end{thebibliography}

\end{document}
